\scanpage{157}
% archive sheet number: 156
% L1–6
[On pourrait dégager
\[
\begin{aligned}
&\text{Soit }s_0\in S,\ D\text{ un rayon. Alors si un \deleted{\unsure{pt}} }t_0\in D\text{ est}\\
&A\text{-adjacent à }s_0,\text{ tous les }t\in D\text{ le sont. L'application}\\
&\qquad D\longrightarrow\Pi\\
&\qquad t\longmapsto\tau_A(\{s_0,t\})\\
&\text{est \uline{bijective}.}\qquad]
\end{aligned}\tag*{(78)}
\]
% L7–10
Ainsi $g\in G_0$ fixe $s_0$ et $t_0$, qui sont adjacents, donc si on avait $g\ne\mathrm{id}$, on aurait $g=\tau_\infty^{\{s_0,t_0\}}$, or $\chi(\tau_\infty)=-1$ \deleted{et} tandis que \deleted{\unsure{$g$}} l'on a supposé $g\in G_0^+$ i.e. $\chi(g)=1$, donc on doit avoir $g=1$.

% L11–13
Prouvons enfin $T\xrightarrow{\chi\mid T}(\mathbb Z/n\mathbb Z)^*=Z$ surjectif. Soit $\lambda^{\unsure{*|\cdot}}\in(\mathbb Z/n\mathbb Z)^*$ \inserted{et soit $\dot\lambda$ son image dans $Z=(\mathbb Z/n\mathbb Z)^*/\pm1$,} alors $\exists\, g\in G_{s_0}$ tel que l'on ait
\[
g\,t_0=\dot\lambda.t_0\tag*{\deleted{(\unsure{74})}\ \unsure{(\ast)|\ast|\text{ou}}}
\]
% L14–17
puisque $G_{s_0}$ est transitif sur $\mathrm{Et}_F(s_0)$, et \unsure{cel|cela} implique que $g\in T$ (75). \deleted{\unsure{A}} La relation $(\ast)$ signifie aussi
\[
\chi'(g)=\dot\lambda\ ,\quad\text{i.e.}\quad\chi(g)\in\pm\lambda\,,
\]
% L18–22
mais d'autre part il y a ambiguïté dans le choix de $g$, via la possibilité
\[
g\longmapsto \deleted{\unsure{\tau}}g\circ\tau_\infty,\qquad\text{donc}\quad\chi(g)\longmapsto\chi(g)\chi(\tau_\infty)
\]
or
\[
\chi(\tau_\infty)=-1
\]
donc en corrigeant au besoin $g$, on gagne !

% L23–27
\uline{Remarque.} Cette façon de trouver une décomposition de $G_{s_0}$ en produit $1/2$ direct $G^+_{0s_0}.\deleted{\unsure{Z}}T_{\unsure{\prime|\,}}$, où $T\subset G_{s_0}$ est tel que $T\longrightarrow\widetilde Z=G/G_0^+$, est ``formelle'' à partir des seuls axiomes $T_0\,1,\ T_0\,4$, et l'axiome de régularité et d'orientabilité $(T_0\,8)$,
